\DTMsavedate{mydate}{2022-06-30}
\maketitle{Engineering Management}{Electronic Engineering}{\DTMusedate{mydate}}

% ----------------------------------------------------- %
% COURSE INFO
\subsection*{Course Information}
\vspace{0ex}%
\noindent\parbox{0.5\textwidth}{%
    \noindent\begin{tabular}{@{}ll}
                 \textsf{Course:}          & 	Engineering Management   \\
                 \textsf{Course Code:}         & 	TEE 5101    \\
                 \textsf{Credit Hours:}    & 	10
    \end{tabular}}
\vspace{2ex}\hrule\vspace{2ex}

% ----------------------------------------------------- %
% INSTRUCTOR INFO
\subsection*{Lecturer Information}
\noindent\parbox{0.5\textwidth}{%
    \noindent\begin{tabular}{@{}ll}
                 \textsf{Name:}         &    Bhekisisa Nyoni         \\
%\textsf{Phone:}        &    \phone          \\
\textsf{Email:}        &    \href{mailto:bhekisisa.dube@nust.ac.zw}{bhekisisa.dube@nust.ac.zw} \\
\textsf{Qualifications:}        &  MPhil in Electronic Engineering (NUST), BEng (Hons) in Electronic Engineering \\
 &                                 (NUST), PGD in Project Management (PMZ), PGD in Higher Education \\
 &                                 (NUST), OSHEMAC (NSSA)

    \end{tabular}}
\vspace{2ex}\hrule\vspace{2ex}


% ----------------------------------------------------- %
% COURSE SYNOPSIS
\subsection*{Course Synopsis}
This module is based on “Management by Engineers” by D. Johnson through group discussion and talks by external speakers.
Hence it is centred on industrial organizations, reviews and performance measures, planning and managing change, development and motivating groups, leaderships and communication, financial management, business environment, companies and basic accounts.
\vspace{2ex} \hrule\vspace{2ex}

% ----------------------------------------------------- %
% COURSE TOPICS
\subsection*{Topics}
\begin{outline}
    \1 Fundamentals of Management
    \begin{itemize}
        \item Engineering
        \item Management
        \item Evolution of Management Thought
    \end{itemize}
    \1 The Organisation
    \1 The Business Environment
    \1 Managing People
    \1 Managing Change
    \1 Total Quality Management
    \1 Financial Management
    \1 Cost
    \begin{itemize}
        \item Break-even analysis
        \item Costing Techniques
    \end{itemize}
\end{outline}
\vspace{2ex} \hrule\vspace{2ex}

% ----------------------------------------------------- %
% Students Assenssment
\subsection*{Assessment of Students}
\begin{outline}
    \1 Continuous assessment will consist of in-class presentations, written assignments and written in-class tests.
    \1 There shall be a final examination at the end of the semester.
\end{outline}
\vspace{2ex}\hrule\vspace{2ex}

% ----------------------------------------------------- %
% Grade Evaluation
\subsection*{Course Evaluation and Grading Policies}
Grades are based on:
Continuous assessment (\qty{25}{\percent}),
Final Exam (\qty{75}{\percent})
\vspace{2ex}\hrule\vspace{2ex}

% ----------------------------------------------------- %
% EQUIPMENT
% https://sites.udel.edu/computing-purchases/personal-specs/
% \subsection*{Essential Equipment}

% \vspace{2ex}\hrule\vspace{2ex}

% Policies and Other information
\subsection*{Policies and Other Information}
% Conduct
\subsection*{Key Text}

\begin{itemize}
    \item Management by R. L. Daft
    \item Modern Project Management by R. C. Mishra, T. Soota
    \item Change Management by B. Burnes
\end{itemize}
\vspace{2ex}\hrule
\vspace{2ex}\hrule\vspace{2ex}